% Adaptado de Semana 06 Clase 02 (2025-A).
\documentclass[compacto,10pt]{pubciencias-notas}
\usepackage{comandosEPN}
\autor{Departamento de Formación Básica}
\materia{Álgebra Lineal}
\nota{Semana 05 Clase 01}
\fecha{Semestre 2026-B}
\logouno[1.1cm]{Logo_EPN.eps}
\logodos[1.7cm]{Logo_DFB.eps}
\definecolor{colordef}{cmyk}{0.81,0.62,0.00,0.22}

\begin{document}
\encabezado
\vspace*{-8mm}
\setcounter{section}{2}
\section{Determinantes}
\setcounter{subsection}{3}

En las semanas anteriores relacionamos la solución de un sistema con el rango
y la matriz inversa. El determinante completa estas conexiones: para una matriz
cuadrada, los siguientes criterios expresan una misma condición.

\begin{teo}[Equivalencias de invertibilidad]\label{sem05cl01:equivalencias}
    Sea $A\in\Mat[\R]{n}{n}$. Son equivalentes:
    \begin{enumerate}
        \item $A$ es invertible o no singular;
        \item $Ax=0$ tiene únicamente la solución trivial;
        \item $A$ es equivalente por filas a $I_n$;
        \item $\rang(A)=n$;
        \item $Ax=b$ tiene una única solución para cada $b\in\R^n$;
        \item $\det(A)\ne0$.
    \end{enumerate}
\end{teo}

\begin{ejem}[Aplicación: brazo robótico]\label{sem05cl01:robot-modelo}
    Retomemos el brazo robótico del ejercicio 5 de la semana 02. Sus dos mecanismos producen
    desplazamientos $u=(2,1)^{\top}$ y $v=(1,3)^{\top}$.
    Las variables $x_1$ y $x_2$ representan, respectivamente, el número de
    activaciones del primer y del segundo mecanismo. Para alcanzar la posición final $b=(8,9)^{\top}$ desde el origen, el sistema asociado es
    \[
        \left\{\begin{aligned}
            2x_1+x_2&=8,\\
            x_1+3x_2&=9.
        \end{aligned}\right.
    \]
    La primera ecuación describe el desplazamiento horizontal y la segunda,
    el vertical. Con $x=(x_1,x_2)^{\top}$, escribimos $Ax=b$, cuya matriz
    de coeficientes es
    \[
        A=\begin{pmatrix}2&1\\1&3\end{pmatrix},
        \qquad \det(A)=2\cdot3-1\cdot1=5\ne0.
    \]
    Por el teorema, $\rang(A)=2$, $A$ tiene inversa y $Ax=0$ solo tiene la
    solución trivial. Además, cada objetivo $b\in\R^2$ determina una única
    solución real de $Ax=b$.
\end{ejem}

Ya sabemos reconocer cuándo existe una solución única. Ahora veremos cómo
obtenerla mediante determinantes.

\subsection{Regla de Cramer}

\begin{teo}[Regla de Cramer]\label{sem05cl01:cramer}
    Sean $A\in\Mat[\R]{n}{n}$ y $b\in\R^n$. Si $\det(A)\ne0$, la única
    solución de $Ax=b$ tiene componentes
    \[
        x_j=\frac{\det(A^{(j)})}{\det(A)},\qquad j=1,\ldots,n,
    \]
    donde $A^{(j)}$ se obtiene reemplazando la columna $j$ de $A$ por $b$,
    sin modificar las demás columnas.
\end{teo}

\noindent\textbf{Procedimiento.}
\begin{enumerate}
    \item Escribir el sistema como $Ax=b$ y comprobar que $\det(A)\ne0$.
    \item Para cada incógnita $x_j$, reemplazar la columna $j$ por $b$.
    \item Calcular el cociente de determinantes y comprobar la solución
          en las ecuaciones originales.
\end{enumerate}

\begin{ejem}[Alcanzar el objetivo del robot]\label{sem05cl01:robot}
    Para alcanzar $b=(8,9)^{\top}$ necesitamos resolver
    \[
        \left\{\begin{aligned}2x_1+x_2&=8,\\x_1+3x_2&=9.\end{aligned}\right.
    \]
    Como $\det(A)=5$, podemos aplicar Cramer. Las matrices de reemplazo son
    \[
        A^{(1)}=\begin{pmatrix}8&1\\9&3\end{pmatrix},
        \qquad A^{(2)}=\begin{pmatrix}2&8\\1&9\end{pmatrix}.
    \]
    Entonces
    \[
        x_1=\frac{\det(A^{(1)})}{\det(A)}=\frac{8\cdot3-1\cdot9}{5}=3,
        \qquad x_2=\frac{\det(A^{(2)})}{\det(A)}=\frac{2\cdot9-8\cdot1}{5}=2.
    \]
    La comprobación es $2(3)+2=8$ y $3+3(2)=9$. Geométricamente,
    \[
        3\begin{pmatrix}2\\1\end{pmatrix}
        +2\begin{pmatrix}1\\3\end{pmatrix}
        =\begin{pmatrix}8\\9\end{pmatrix}.
    \]
    Tres activaciones del primer mecanismo y dos del segundo permiten
    alcanzar el objetivo; en este caso las activaciones son enteras y no negativas.
\end{ejem}

\subsection{¿Qué ocurre si el determinante es cero?}

\begin{advertencia}
    Cramer requiere $\det(A)\ne0$. Si $\det(A)=0$, la fórmula no se aplica.
    Esto no permite afirmar que el sistema carece de solución: puede tener infinitas soluciones o ninguna.
\end{advertencia}

Cramer resulta útil para sistemas pequeños, para sistemas grandes, la eliminación gaussiana evita calcular muchos determinantes y sigue siendo nuestra herramienta principal.

\section{Ejercicios}

\begin{ejer}[Cramer: tres incógnitas]
    Considere
    \[
        \left\{\begin{aligned}
            3x+2y+z&=2,\\2x+2y+z&=2,\\x+y+z&=1.
        \end{aligned}\right.
    \]
    \begin{enumerate}
              \item Identifique la matriz de coeficientes
              $A$ y el vector $b$. Calcule $\det(A)$ y compruebe que vale 1.
              Para simplificar este cálculo, puede aplicar
              $-F_2+F_1\rightarrow F_1$ y luego expandir por cofactores en la nueva primera fila. Esta operación conserva el
              determinante.
        \item Usando $A$ y $b$, aplique la regla de Cramer para determinar $x$, $y$ y $z$. Compruebe la solución en las tres ecuaciones originales.
    \end{enumerate}
\end{ejer}

\begin{ejer}[Aplicación: mezcla]
    Se desea preparar 10 litros de una mezcla al $30\%$ a partir de dos soluciones, una al $20\%$ y otra al $40\%$. Suponga que los volúmenes se suman. Si $x$ e $y$ son los litros de cada solución:
    \begin{enumerate}
            \item Justifique el sistema
              \[
                  \left\{\begin{aligned}
                      x+y&=10,\\
                      0.2x+0.4y&=3.
                  \end{aligned}\right.
              \]
        y resuelva por Cramer. Interprete las unidades de la respuesta.
        \item Si ambas soluciones iniciales fueran al $20\%$, ¿podría
              obtenerse la mezcla solicitada? Relacione su respuesta con
              el determinante y con la concentración final.
    \end{enumerate}
\end{ejer}

\end{document}
